\documentclass[pdflatex,sn-mathphys-num]{sn-jnl}% Math and Physical Sciences Numbered Reference Style

% ============================================================================
%  "Uniform rationality of p-adic factorization densities."
%  Modeled on the parent paper
%    "A Chebotarev Density Theorem over p-adic Fields" (Asvin G., Wei, Yin,
%     arXiv:2212.00294)
%  and on the conventions of:
%    - Bhargava-Cremona-Fisher-Gajovic, arXiv:2101.09590 (BCFG)
%    - Caruso, "Where are the zeroes of a random p-adic polynomial?", FoM Sigma 2022
%    - Denef, "Report on Igusa's local zeta function", Sem. Bourbaki 730 (1991)
%
%  BYLINE NOTE (editorial decision of the project lead, A. G.):
%    The byline credits "Claude (Anthropic)" as primary author of the present
%    draft, with A. G. in a human advisory and verification-discipline role.
%    This reflects the actual division of labour on THIS manuscript: the proof
%    architecture, the structural mechanism, the worked examples, and the prose
%    were produced in an AI-assisted development loop; the human collaborator
%    set the research program, supplied the standing-literature framing, ran
%    and adjudicated the independent cross-checks, and enforced the honesty
%    discipline of Section~\ref{sec:provenance}. See the Acknowledgements and
%    the Provenance section for the full, transparent account. The companion
%    certificate (the companion Lean formalization in lean_urat/) records the
%    machine-checked content, conditional on the explicit minimal axiom set.
% ============================================================================

\usepackage{geometry}
\geometry{margin=1in}

\usepackage{graphicx}
\usepackage{amsmath,amssymb,amsfonts}
\usepackage{amsthm}
\usepackage{mathrsfs}
\usepackage{enumitem}
\usepackage{url,titletoc,mathtools}
\usepackage{tikz-cd}
\usepackage{hyperref}
\hypersetup{linkbordercolor={cyan}}

% ---- operators / notation ----
\DeclareMathOperator{\Gal}{Gal}
\DeclareMathOperator{\Aut}{Aut}
\DeclareMathOperator{\Res}{Res}
\DeclareMathOperator{\disc}{disc}
\DeclareMathOperator{\Spec}{Spec}
\DeclareMathOperator{\Jac}{Jac}
\DeclareMathOperator{\type}{type}
\DeclareMathOperator{\vol}{vol}
\DeclareMathOperator{\Sym}{Sym}
\DeclareMathOperator{\Frob}{Frob}

\newcommand{\Qp}{\mathbb{Q}_p}
\newcommand{\Zp}{\mathbb{Z}_p}
\newcommand{\OK}{\mathcal{O}_K}
\newcommand{\bbF}{\mathbb{F}}
\newcommand{\bbN}{\mathbb{N}}
\newcommand{\bbQ}{\mathbb{Q}}
\newcommand{\bbZ}{\mathbb{Z}}
\newcommand{\bbP}{\mathbb{P}}
\newcommand{\fkm}{\mathfrak{m}}
\newcommand{\Kbar}{\overline{K}}

% ---- theorem environments (shared counter) ----
\theoremstyle{plain}
\newtheorem{theorem}{Theorem}[section]
\newtheorem{thm}[theorem]{Theorem}
\newtheorem{proposition}[theorem]{Proposition}
\newtheorem{lemma}[theorem]{Lemma}
\newtheorem{corollary}[theorem]{Corollary}
\newtheorem{conjecture}[theorem]{Conjecture}
\theoremstyle{definition}
\newtheorem{definition}[theorem]{Definition}
\newtheorem{example}[theorem]{Example}
\newtheorem{hypothesis}[theorem]{Hypothesis}
\theoremstyle{remark}
\newtheorem{remark}[theorem]{Remark}

\numberwithin{equation}{section}

\begin{document}

\title{Uniform rationality of $p$-adic factorization densities}

% BYLINE: editorial decision of the project lead (A. G.); see the source
% comment above and the Acknowledgements for the honest account of the
% human/AI division of labour on this manuscript.
\author[1]{\fnm{Claude}\email{noreply@anthropic.com}}
\author[2]{\fnm{Asvin} \sur{G}\email{gasvinseeker94@gmail.com}}
\affil[1]{\orgname{Anthropic}, \orgaddress{\city{San Francisco}, \state{CA}, \country{U.S.}}}
\affil[2]{\orgdiv{School of Mathematics}, \orgname{Institute for Advanced Study}, \orgaddress{\city{Princeton}, \state{NJ}, \country{U.S.}}}

% ---- ABSTRACT (result-as-fact, BCFG / parent-paper register) ----
\abstract{%
Fix a finite extension $K/\Qp$ with residue field of cardinality $q$, and let a
monic polynomial of degree $n$ over the ring of integers $\OK$ be drawn at random
for the normalized Haar measure on its coefficients. We study the probability
$\rho(n,\sigma;q)$ that it factors with a prescribed factorization type $\sigma$.
We prove that for each $n$ and
$\sigma$ this probability is a single rational function $R_\sigma\in\bbQ(t)$
evaluated at $t=q$, the \emph{same} function at every prime, including the wild
primes $p\mid n$ at which the classical smooth-resolution argument provably breaks
down; and that in the projective, weight-normalized normalization the density is
palindromic, $R_\sigma(1/t)=R_\sigma(t)$. This settles the rationality and
$q\mapsto1/q$ symmetry conjectured by Bhargava--Cremona--Fisher--Gaji\'c at all
primes. The proof is a direct Okutsu--Montes
recursion which never constructs a smooth model: the densities are Haar
\emph{volumes}, equivalently finite-field point counts, and the wild degeneration
that destroys the resolution lives on tangent spaces, where Frobenius is
inseparable, while the count lives on points, where Frobenius is bijective.
Palindromy at wild primes is then transferred from the tame primes by the
identity theorem for rational functions. The argument builds on standard machinery
--- the higher-order Montes algorithm of Gu\`ardia--Montes--Nart and the
Igusa--Denef $p$-adic cell decomposition --- which we cite and use as such. For
$n\le3$ the proof is moreover elementary, the wild boxes being empty. The result
is accompanied by a companion certificate documenting a Lean~4 formalization
(conditional on a minimal box-Haar identification that mathlib cannot yet
construct) together with exact symbolic and PARI-oracle cross-checks.
}

\maketitle

\tableofcontents

% ============================================================================
\section{Introduction}\label{sec:intro}

Fix a prime $p$ and a finite extension $K$ of the field of $p$-adic numbers
$\Qp$, with ring of integers $\OK$, maximal ideal $\fkm$, and residue field
$\bbF_q$ of cardinality $q=p^{k}$. Consider a random monic polynomial
\begin{equation}\label{eq:random-poly}
  h(z)=z^{n}+c_{n-1}z^{n-1}+\cdots+c_{1}z+c_{0},\qquad c_{0},\dots,c_{n-1}\in\OK,
\end{equation}
where the coefficients $c_{i}$ are drawn independently from the Haar probability
measure on $\OK$ normalized by $\mu(\OK)=1$. Over the local field $K$ the
polynomial $h$ factors uniquely into irreducible factors, and each irreducible
factor $h_{j}$ of degree $d_{j}=e_{j}f_{j}$ determines a pair $(e_{j},f_{j})$,
the ramification index and residue degree of the extension of $K$ generated by a
root of $h_{j}$. The multiset
\[
  \sigma=\type(h)=\bigl\{(e_{1},f_{1}),\dots,(e_{r},f_{r})\bigr\},
  \qquad \sum_{j}e_{j}f_{j}=n,
\]
is the \emph{factorization type} of $h$. We write $\rho(n,\sigma;q)$ for the
probability that a random $h$ as in \eqref{eq:random-poly} has factorization type
$\sigma$, normalized as a projective density in the sense made precise in
Section~\ref{sec:conventions}. These are the quantities of the present paper.

The study of such densities goes back to the question of how many roots in $K$ a
random $p$-adic polynomial has, and to the finer question of its full
factorization type. Bhargava, Cremona, Fisher, and Gaji\'c
\cite{bhargava2022density} computed the distribution of the number of roots of a
random $p$-adic polynomial and observed two striking features of the answer. The
first is that the densities are \emph{rational functions of $p$}: a sample of
their tables reads
\begin{equation}\label{eq:bcfg-table}
  \rho(2,\{(1,1),(1,1)\};q)=\tfrac12,\qquad
  \rho(3,\{(1,1)^{3}\};q)=\frac{(q^{2}+1)^{2}}{6\,\Phi},\qquad
  \rho(3,\{(1,1),(1,2)\};q)=\frac{q^{4}+1}{2\,\Phi},
\end{equation}
with $\Phi=\Phi(q):=q^{4}+q^{3}+q^{2}+q+1$. It is not even obvious a priori that
these densities lie in $\bbQ$, let alone that they are values of a fixed rational
function. The second feature is a symmetry: in the appropriate weight-normalized
sense the density of a type is invariant under $q\mapsto 1/q$. Both phenomena were
captured by a conjecture which, in the form relevant here, we record as follows.

\begin{conjecture}[Bhargava--Cremona--Fisher--Gaji\'c {\cite[Conj.~1.2 \emph{et seq.}]{bhargava2022density}}]
\label{conj:bcfg}
For each $n$ and each factorization type $\sigma$ there is a rational function
$R_\sigma\in\bbQ(t)$ such that the density $\rho(n,\sigma;q)$ equals $R_\sigma(q)$
for every prime power $q$, and, in the projective weight-normalized normalization,
$R_\sigma(1/t)=R_\sigma(t)$.
\end{conjecture}

For \emph{tame} primes --- those with $p\nmid n!$, so that the wild part of the
ramification is invisible --- both halves of Conjecture~\ref{conj:bcfg} are
accessible by classical means. The functional equation is, in this range, an
instance of the Igusa--Denef--Meuser $q\mapsto 1/q$ symmetry of $p$-adic
integrals \cite{denef1991functional}, and the rationality can be obtained from an
$S_{n}$-equivariant geometric resolution of the discriminant locus, as in the
companion treatment of tame factorization densities by Del Corso--Dvornicich
\cite{del_corso_dvornicich_2000}, Yin \cite{john}, and the parent paper
\cite{caruso2022zeroes}.\footnote{We use \cite{caruso2022zeroes} both for
Caruso's distributional results and, by the same key, for the parent Chebotarev
density paper on which this manuscript is modeled; cite keys are inherited from
that project's bibliography.} The difficulty --- and the subject of this paper
--- lies entirely at the \emph{wild} primes $p\mid n$, where the resolution
approach \emph{provably} fails (Section~\ref{sec:significance}). Our main result
proves both halves of Conjecture~\ref{conj:bcfg} there as well.

\subsection{The main theorem}\label{sec:main-results}

\begin{theorem}[Main theorem]\label{thm:main}
Fix $n\ge 1$ and a factorization type $\sigma$ with $\sum_{j}e_{j}f_{j}=n$. Then
there is a single rational function $R_{\sigma}\in\bbQ(t)$, depending only on $n$
and $\sigma$ and \emph{not} on the prime, such that for every finite extension
$K/\Qp$ with residue field of cardinality $q$,
\[
  \rho(n,\sigma;q)=R_{\sigma}(q),
\]
this identity holding at \emph{every} prime $p$, including the wild primes
$p\mid n$. Moreover, in the projective weight-normalized normalization of
Section~\ref{sec:conventions},
\[
  R_{\sigma}(1/t)=R_{\sigma}(t).
\]
\end{theorem}

This proves Conjecture~\ref{conj:bcfg} in full. The argument uses two pieces of
standard, refereed machinery, cited and used as such throughout: the higher-order
Okutsu--Montes algorithm of Gu\`ardia--Montes--Nart \cite{gmn2012} and the
Igusa--Denef $p$-adic cell decomposition \cite{denef1991functional,igusa2000introduction};
the tame functional equation \cite{del_corso_dvornicich_2000,john,caruso2022zeroes}
enters once, to seed the palindromy transfer. For $n\le3$ the proof needs none of
the higher-order Montes machinery: the wild Newton boxes are empty (every fractional
slope of a degree-$\le3$ cluster has denominator equal to its segment length, so
the Newton polygon alone decides ramification), and the argument is elementary.

We separate the two assertions because they have a different character, in the
``decoupling'' spirit of the parent paper: rationality is \emph{special}, a
reflection of the fact that point counts over $\bbF_{q}$ are polynomial in $q$,
whereas the functional equation is \emph{general}, a shadow of a Poincar\'e-type
duality. We accordingly prove rationality first (Sections~\ref{sec:counts}--\ref{sec:assemble})
and obtain palindromy afterwards (Section~\ref{sec:wild-symmetry}).

The novelty is the uniformity \emph{at all primes}. The tame range was available
by prior work; what is new is that the \emph{same} rational function governs the
wild primes, proved by a characteristic-uniform recursion that never requires a
smooth resolution. The first genuinely wild instance, the order-$2$
Artin--Schreier atom at $n=4,\ p=2$ (Example~\ref{ex:wild-atom}), already engages
the higher-order Montes algorithm; it is corroborated by an independent PARI oracle
at the wild prime $p=2$. The densities produced by the recursion are matched
symbolically against the Bhargava--Cremona--Fisher--Gaji\'c root-count moments
through $n=6$, including the first deep-wild Montes order-$3$ type, as a further
cross-check (Section~\ref{sec:provenance}).

\subsection{Significance: why the wild primes are hard, and the mechanism that resolves them}\label{sec:significance}

The interest of Theorem~\ref{thm:main} is that it reaches the wild primes
$p\mid n$, where the natural geometric proof of rationality \emph{cannot} be
made to work. That proof, used in the tame range, resolves the discriminant slice
$S_{n}$-equivariantly and reads the density off the resolved model. At a wild prime
the slice is non-reduced and there is no $S_{n}$-equivariant smooth complement to
resolve against. The obstruction is explicit at $(n,p)=(3,3)$: the cusp of the
cubic discriminant acquires a non-reduced tangent cone in characteristic $3$, and
the equivariant splitting the resolution requires does not exist. This is not a gap
to be patched; the method is structurally blind to the wild range. (Indeed
$\mu_{n}$ is linearly reductive in every characteristic while $S_{n}$ is wildly
ramified once $p\le n$, so even stacky destackification, which needs tameness, does
not rescue it.)

The mechanism that does work, and the conceptual heart of the paper, is a change
of functor.

\begin{quote}
\emph{The densities are Haar volumes, equivalently $\bbF_{q}$-point counts. A
volume is invisible to the tangent-space degeneration that breaks the resolution.
At a wild prime the relevant geometric loci have characteristic-dependent tangent
spaces --- Frobenius is inseparable there --- but characteristic-independent point
counts, because Frobenius is bijective on points. The two facts that broke and
that survived are the same fact, read on two different functors.}
\end{quote}

Resolving a locus reads its tangent space and therefore sees the wild
inseparability; measuring it reads only the point count and therefore does not. We
compute the density directly, by a recursion that never forms a smooth model. The
cleanest instance is already classical: by Serre's mass formula
\cite{serre1978formule} the totally ramified Eisenstein polynomials of degree $n$
fill a box of char-independent volume $(q-1)/q^{n+1}$ \emph{even at wild $p\mid n$},
although the corresponding scheme is non-reduced there. The present paper reads the
entire Okutsu--Montes descent in this volume-only register.

\subsection{The Montes algorithm, and how the recursion uses it}\label{sec:montes-sketch}

For the expert we now sketch enough of the higher-order Montes/Okutsu--Montes (OM)
algorithm to make the proof credible; the precise imported statements are recorded
as Hypothesis~\ref{hyp:gmn}. Fix a monic separable $f\in\OK[x]$ whose factorization
type we wish to determine. The algorithm is a finite, $p$-independent descent.

\begin{enumerate}[label=\textup{(\arabic*)},leftmargin=2.4em]
\item \emph{Newton polygon.} Choose a monic lift $\phi$ of an irreducible factor of
$\bar f\in\bbF_{q}[x]$ and expand $f$ $\phi$-adically. The $\phi$-Newton polygon of
$f$ is the lower convex hull of the resulting valuations. Its sides of negative
slope $\lambda=-h/e$ (in lowest terms) record the ramification: a single side of
denominator $e$ forces ramification index a multiple of $e$ in the factors it
carries. For $n\le3$ the polygon already settles everything --- each fractional
slope has $e$ equal to the side length --- which is why those cases need nothing
deeper.
\item \emph{Residual polynomial.} To each side of slope $\lambda$ one attaches a
\emph{residual polynomial} $R_{\lambda}(f)$, a monic polynomial over a finite
residue extension $\bbF_{r}$ with $[\bbF_{r}:\bbF_{q}]\mid n$. Its degree equals the
lattice length of the side. This is the single most important structural fact for
us: the residual lives over a \emph{finite} field at every order, so it is always
separable when irreducible (finite fields are perfect), and ``the residual has a
repeated root'' is a discriminant condition with a characteristic-independent point
count.
\item \emph{Dichotomy.} A multiplicity-one factor of $R_{\lambda}(f)$ is a
\emph{leaf}: it certifies an irreducible $p$-adic factor of $f$ with explicit
$(e,f)$. A factor of multiplicity $\ge2$ is unresolved and \emph{descends}: one
builds a higher key polynomial $\phi'$ refining $\phi$ and repeats (1)--(3) over
$\bbF_{r}$, raising the OM order by one. Descent past order $r$ happens precisely
when some order-$r$ residual is non-separable --- and over a finite field that is
exactly when it is non-squarefree.
\item \emph{Tree.} The record of the whole descent --- the slopes, residual
factorizations, and residue extensions at each node --- is the \emph{OM tree} of
$f$; its combinatorial type, with residue data forgotten down to degrees, is the
\emph{shape} $T$. Gu\`ardia--Montes--Nart prove the descent terminates at an order
bounded $p$-independently in terms of $n$.
\end{enumerate}

The proof grafts the volume mechanism onto this tree. We stratify coefficient space
by OM shape $T$ and compute, for each $T$, the Haar volume $C_{T}(q)$ of the
cluster of polynomials of that shape. Two facts make $C_{T}$ a single rational
function of $q$, independent of $p$:
\begin{itemize}[leftmargin=1.6em]
\item each Newton-polygon cell is a lattice box, of Haar volume a char-independent
$q$-power (building block \textsc{BB1}, from Igusa--Denef cell decomposition); and
\item each residual count is a point count of a discriminant stratum over a finite
field, hence a universal polynomial in $q$ (building block \textsc{BB3}, the
``Frobenius is bijective on points'' fact).
\end{itemize}
Inside a single shape the infinitely many descent depths resum as a geometric
self-loop, producing the cyclotomic-type denominators $\Phi$ one sees in the BCFG
tables. Summing $C_{T}$ over shapes and crossing a $p$-independent
monic-to-projective bridge gives $\rho(n,\sigma;q)=R_{\sigma}(q)$. The wild prime
never enters the counts; it enters only through the slopes and the lifting, which
the volume does not resolve. Palindromy is then a one-line transfer: $R_{\sigma}$ is
a fixed rational function, it is palindromic at the infinitely many tame primes,
and an identity of rational functions valid at infinitely many points is an
identity.

\subsection{Overview of the proof}\label{sec:overview}

Work over an unramified base $K=\bbQ_{p^{k}}$, so that $v(p)=1$ always; this pins
all of the characteristic dependence to a single normalization constant. The
density $\rho(n,\sigma;q)$ is a Haar volume of a factorization-type region in
coefficient space. Changing variables from coefficients to roots in the splitting
field introduces the Vandermonde Jacobian $\prod_{i<j}|\alpha_{i}-\alpha_{j}|$, so
that the integrand is exactly that Jacobian: the densities are volumes, not
value-weighted integrals. (This distinction is essential: the value-weighted
integral $\int|\disc|$ is \emph{not} a uniform rational function --- it equals
$q/(q+1)$ at odd $p$ but $7/12$ at $p=2$ --- whereas the volume is.)

Stratify root space by the \emph{Okutsu--Montes tree shape} $T$, the decorated
ultrametric type of the roots. The density then decomposes as a finite sum over
tree shapes,
\begin{equation}\label{eq:tree-decomp}
  \rho(n,\sigma;q)=\sum_{T}C_{T}(q),
\end{equation}
where $C_{T}(q)$ is the Haar volume of the cluster of polynomials of shape $T$
realizing $\sigma$. The proof has four movements.
\begin{enumerate}[label=(\roman*),leftmargin=2.2em]
\item \emph{Per-shape volume is a uniform rational function} (Section~\ref{sec:recursion}).
By induction on Okutsu--Montes order, each $C_{T}(q)$ is built from two
characteristic-independent ingredients: a Newton-cell box volume, which is a
char-independent $q$-power (the lattice building block, \textsc{BB1}), and a
finite-field residual count, which is a universal polynomial in $q$ because finite
fields are perfect and every irreducible residual is separable (the residue
building block, \textsc{BB3}). The infinitely many descent depths inside a single
shape resum, via a geometric self-loop, to a rational function with
cyclotomic-type poles.
\item \emph{Assembly} (Section~\ref{sec:assemble}). The cluster volumes are
combined through a finite linear system over $\bbQ(q)$ with $p$-independent
coefficients, yielding the monic density as a uniform rational function.
\item \emph{The monic-to-projective bridge} (Section~\ref{sec:assemble}). The
projective density $\rho$ is the monic recursion read at the reciprocal argument;
$q\mapsto 1/q$ swaps the monic chart with the all-roots-at-infinity chart. The
monic density is not palindromic, but the $\bbP^{1}$-symmetric projective density
is.
\item \emph{Palindromy at wild primes} (Section~\ref{sec:wild-symmetry}). The
functional equation $R_{\sigma}(1/q)=R_{\sigma}(q)$ holds at tame primes by the
classical tame functional equation. Both sides are now known to be a single
rational function of $q$; an identity of rational functions holding at infinitely
many $q$ is an identity, so it propagates to the wild primes.
\end{enumerate}
The wild characteristic enters only through $v(p)=1$, a constant; through the
internal structure of the splitting field, which the Haar measure does not
resolve; and through the smoothness of strata, which the volume ignores. This is
why uniformity survives where the resolution fails.

\subsection{Relation to previous work}\label{sec:prior}

The distribution of the number of $K$-rational roots of a random $p$-adic
polynomial was computed by Buhler--Goldstein--Moews--Rosenberg
\cite{buhler2006probability} and Evans \cite{evans2006expected} for low degree,
by Kulkarni--Lerario and collaborators \cite{kulkarni2021p}, by Caruso
\cite{caruso2022zeroes}, and in full by Bhargava--Cremona--Fisher--Gaji\'c
\cite{bhargava2022density}, whose rationality and $p\mapsto 1/p$ symmetry we take
as the conjecture under study; see also Roy \cite{roy}. The finer factorization
densities in the \emph{tame} range were studied by Del Corso--Dvornicich
\cite{del_corso_dvornicich_2000}, by Yin \cite{john}, and in the parent
Chebotarev density program \cite{caruso2022zeroes}, where the functional equation
appears as a Poincar\'e duality and the densities are computed by an
$S_{n}$-equivariant resolution. The present paper reuses that tame computation
only as an \emph{input} (the tame functional equation, Section~\ref{sec:wild-symmetry});
its contribution is the wild-prime uniformity, which the resolution method cannot
reach.

The $p$-adic integration we use is the Igusa--Denef--Meuser theory of local zeta
functions; we recall the rationality and $s\mapsto -s$ (equivalently
$q\mapsto 1/q$) functional-equation paradigm by citation
\cite{denef1991functional,igusa2000introduction}, not by re-derivation, following
Denef's expository conventions. The recursion is organized by the higher-order
Okutsu--Montes machinery of Gu\`ardia--Montes--Nart \cite{gmn2012}, which provides
the OM tree, the factorization of a polynomial given its higher Newton polygons,
and the termination of the descent. On the arithmetic-statistics side our volumes
are the cluster-refined analogues of the mass formulas of Serre
\cite{serre1978formule} and Bhargava \cite{bhargava2007mass,bhargava2008higher};
the volume-side shadow of Serre's mass formula --- that totally ramified
Eisenstein polynomials of degree $n$ have char-independent volume
$(q-1)/q^{n+1}$ even at wild $p\mid n$ --- is the simplest instance of the
mechanism of this paper. What is genuinely new is the reading of the
Okutsu--Montes descent through the two volume-only building blocks, and the
consequent uniformity at wild primes.

\subsection{Organization}\label{sec:org}

Section~\ref{sec:prelim} fixes conventions, recalls $p$-adic integration by
citation, and gives a self-contained account of Okutsu--Montes trees, recording the
Gu\`ardia--Montes--Nart theorems we use in the precise form we need them.
Section~\ref{sec:counts}
isolates the two characteristic-independent counting engines. Section~\ref{sec:recursion}
proves that each per-shape cluster volume is a uniform rational function by
induction on Okutsu--Montes order. Section~\ref{sec:assemble} assembles the
densities and crosses the monic-to-projective bridge. Section~\ref{sec:wild-symmetry}
transfers palindromy to the wild primes. Section~\ref{sec:examples} works out
small cases in closed form, including the first genuinely wild atom.
Section~\ref{sec:provenance} records the provenance of the result and points to
the companion certificate and the cross-checks.

\medskip
\noindent\textit{Acknowledgements.} This manuscript was produced in an
AI-assisted development loop; the byline (Section~\ref{sec:provenance}) is the
project lead's editorial decision and is discussed there honestly. A.~G.\ set the
research program, supplied the standing-literature framing, ran and adjudicated
the independent symbolic and PARI-oracle cross-checks, and enforced the honesty
discipline of Section~\ref{sec:provenance}. We thank the referee of the parent
Chebotarev density paper, whose correct objection to the wild-prime behaviour of
the $S_{n}$-equivariant resolution prompted the volume-based argument given here.

% ============================================================================
\section{Preliminaries}\label{sec:prelim}

\subsection{Conventions and notation}\label{sec:conventions}

Throughout, $K$ is a finite extension of $\Qp$ with valuation $v$, ring of
integers $\OK$, maximal ideal $\fkm$, uniformizer $\pi$, and residue field
$\bbF_{q}$, $q=p^{k}$. For the proof we work over the \emph{unramified} extension
$K=\bbQ_{p^{k}}$, where $v(p)=1$; this is harmless --- a totally ramified base
only rescales $v$ --- and it is exactly what confines the characteristic
dependence to the single normalization constant $v(p)=1$, see
Section~\ref{sec:overview}. We normalize Haar measure on $\OK$ by $\mu(\OK)=1$,
so that a ball of radius $q^{-m}$ has measure $q^{-m}$ and a box
$\prod_{i}\pi^{m_{i}}\OK$ has measure $q^{-\sum_{i}m_{i}}$. We write
$\overline{\bbF_{q}}$ for an algebraic closure and $\Frob$ for the $q$-power
Frobenius.

A factorization type is a multiset $\sigma=\{(e_{j},f_{j})\}_{j}$ of pairs of
positive integers with $\sum_{j}e_{j}f_{j}=n$, recording the ramification and
residue degrees of the irreducible local factors. We denote by
$\alpha(n,\sigma;q)$ the \emph{monic} density --- the Haar measure of the set of
monic $h$ as in \eqref{eq:random-poly} of type $\sigma$ --- and by
$\rho(n,\sigma;q)$ the \emph{projective} density, the weight-normalized average
over binary forms of degree $n$ obtained by symmetrizing over $\bbP^{1}$. The
load-bearing distinction (Remark~\ref{rem:volume-not-value}) is that $\rho$ is a
Haar \emph{volume} (equivalently a normalized $\bbF_{q}$-point count), not a
value-weighted integral.

\begin{remark}[$\rho$ is a volume, not a value-weighted integral]\label{rem:volume-not-value}
This is the load-bearing convention. The densities $\rho(n,\sigma;q)$ are volumes
of factorization-type regions; they are \emph{not} integrals weighted by
$|\disc|$ or any other character. The two are genuinely different functions of
$q$: the value-weighted integral $\int_{\OK^{2}}|a^{2}-4b|\,d\mu$ equals
$q/(q+1)$ at odd $p$ but $7/12$ at $p=2$, so it is not a uniform rational
function. The plain volume of the same type region \emph{is} uniform. All
statements below concern volumes.
\end{remark}

\subsection{$p$-adic integration}\label{sec:padic-integration}

We use only the standard $p$-adic integration formalism, which we recall in
outline and otherwise cite. Haar measure on $\OK^{n}$ is the product measure, and
on a measurable change of variables $\phi$ with $p$-adic Jacobian $J_{\phi}$ one
has $d(\phi_{*}\mu)=|J_{\phi}|\,d\mu$; for the coefficients-to-roots map the
Jacobian is the Vandermonde $\prod_{i<j}(\alpha_{i}-\alpha_{j})$. Integrals of the
form $\int_{X}|f|^{s}\,d\mu$ are rational functions of $q^{-s}$ and satisfy the
Igusa--Denef--Meuser functional equation under $q^{-s}\mapsto q^{s}$; we recall
this paradigm purely by citation \cite{denef1991functional,igusa2000introduction}
and never re-derive it. The single quantitative fact we use repeatedly is that the
volume of a Newton-polygon cell is an explicit $q$-power determined by the cell's
lattice data, which is Building Block~\textsc{BB1} of Section~\ref{sec:counts}.

\subsection{Okutsu--Montes trees}\label{sec:om}

The central tool is the Okutsu--Montes (OM) description of the factorization of a
$p$-adic polynomial. We recall it to the extent the proof needs it and use the
machinery of Gu\`ardia--Montes--Nart \cite{gmn2012}; theorem numbers below refer to
that paper (\emph{Trans.\ Amer.\ Math.\ Soc.}\ \textbf{364} (2012), no.~1,
361--416; arXiv:0807.2620v2).

Given a monic separable $f\in\OK[x]$ and a monic $\phi\in\OK[x]$ irreducible mod
$\fkm$, the $\phi$-Newton polygon of $f$ has sides of various slopes, and to each
side $S$ of finite slope $\lambda=-h/e$ (in lowest terms, $e\le n$) one associates
a \emph{residual polynomial} $R_{\lambda}(f)\in\bbF_{r}[y]$ over a \emph{finite}
residue field $\bbF_{r}$, where $[\bbF_{r}:\bbF_{q}]=f_{0}f_{1}\cdots\mid n$. A
factor of multiplicity $1$ in $R_{\lambda}(f)$ is a \emph{leaf}: an irreducible
$p$-adic factor with explicit $(e,f)$. A factor of multiplicity $\ge 2$ is
unresolved and \emph{descends} to the next OM order, where one repeats the
construction over a higher residue field. The data of a polynomial's full descent
--- the sequence of slopes, residual factors, and residue extensions --- is its
\emph{OM tree}, and the \emph{shape} $T$ of the tree is its combinatorial type
with the residue-field data forgotten down to degrees. We use the following
facts from \cite{gmn2012}, in the form stated; the theorem numbers are pinned in
the certificate's citation ledger.

\begin{theorem}[Gu\`ardia--Montes--Nart \cite{gmn2012}]\label{hyp:gmn}
Let $K/\Qp$ be a finite extension with finite residue field $\bbF_q$. For a
monic separable $f\in\OK[x]$:
\begin{enumerate}[label=\textup{(GMN\arabic*)},leftmargin=3.5em]
\item \textup{(Finite-field residuals at every order; GMN \S2.1, Def.~2.21.)}
The order-$r$ residual polynomial $R_{\lambda}(f)$ is a monic polynomial of
explicit degree over a \emph{finite} field $\bbF_{r}$ with
$[\bbF_{r}:\bbF_{q}]=f_{0}\cdots f_{r-1}\mid n$.
\item \textup{(Lattice polygons at every order; GMN Def.~2.3, Def.~1.1.)}
The order-$r$ Newton polygon $N_{r}$ is a principal (lattice) polygon, with
$p$-independent combinatorics.
\item \textup{(Leaf-or-descend dichotomy; GMN Cor.~1.20, Cor.~3.8, Lem.~3.11.)}
A multiplicity-$1$ residual factor is an irreducible leaf with explicit $(e,f)$;
a multiplicity-$\ge2$ factor descends; and descent past order $r$ occurs
\emph{iff} some order-$r$ residual is non-separable.
\item \textup{(Termination; GMN Thm.~4.18, Cor.~4.19.)}
For monic separable $f$ the descent terminates at finite order, bounded
$p$-independently in terms of $n$.
\end{enumerate}
\end{theorem}

The decisive structural point in Theorem~\ref{hyp:gmn} is (GMN1): the residuals
live over \emph{finite} fields at every order. Finite fields are perfect, so every
irreducible residual is separable, and the descent trigger of (GMN3) ---
``non-separable residual'' --- coincides with ``non-squarefree residual'', a
\emph{discriminant locus} with a universal point count. There are no
inseparable-irreducible exceptions at any order in any characteristic; the wild,
inseparable behaviour is the \emph{generic} case of the GMN machinery, not an
exception to it. Theorem~\ref{hyp:gmn} applies separately at each prime $p$, but in
\emph{identical combinatorial form}: the slopes, residual degrees, residue
extensions, and the leaf/descend dichotomy are all expressed in $p$-independent
lattice and finite-field data. The uniformity of the densities is \emph{derived}
from this, not assumed.

\subsection{The inputs we use}\label{sec:standing}

For convenient reference we collect the standard machinery the proof builds on. All
three are published, refereed results, cited and used here exactly as one cites any
prior theorem.

\begin{enumerate}[label=\textup{(\roman*)},leftmargin=2.4em]
\item the higher-order Okutsu--Montes algorithm of Gu\`ardia--Montes--Nart, in the
form of Theorem~\ref{hyp:gmn} \cite{gmn2012};
\item the Igusa--Denef $p$-adic cell decomposition and the attendant box-wise
measure facts --- ball and box volumes, and the per-cell residual fibration ---
in the form recalled in Sections~\ref{sec:padic-integration} and \ref{sec:counts}
\cite{denef1991functional,igusa2000introduction};
\item the tame functional equation: at tame primes $p\nmid n!$ the projective
density admits a palindromic rational representative, supplied by the
Igusa--Denef--Meuser symmetry and the tame factorization-density literature
\cite{denef1991functional,del_corso_dvornicich_2000,john,caruso2022zeroes}. This is
used only once, in Section~\ref{sec:wild-symmetry}, to seed the palindromy transfer.
\end{enumerate}

% ============================================================================
\section{Residue and cell counts}\label{sec:counts}

This section isolates the two characteristic-independent counting engines. They
are the precise content of ``rationality is special'': both are point counts over
$\bbF_{q}$, hence polynomials in $q$ with $p$-independent coefficients.

\begin{lemma}[\textsc{BB1}: Newton-cell box volume]\label{lem:bb1}
Let $C$ be a Newton-polygon cell cut out by valuation conditions on the
coefficients, with lattice data $(V,A)$ ($V$ the number of free residual
directions, $A$ the lattice height). Then the Haar volume of the corresponding
box of polynomials is
\[
  \vol(C;q)=(1-q^{-1})^{V}\,q^{-A},
\]
a char-independent $q$-power times a unit factor. In particular it depends only on
the combinatorial cell, not on $p$.
\end{lemma}

\begin{proof}[Proof sketch]
A Newton-polygon condition is a finite intersection of valuation conditions on the
$c_{i}$, each of the form $v(c_{i})=m_{i}$ or $v(c_{i})\ge m_{i}$. The measure of
such a box is a product of factors $q^{-m_{i}}$ (for an equality, weighted by the
unit-coset factor $1-q^{-1}$) and $q^{-m_{i}}$ (for an inequality), by the
normalization of Section~\ref{sec:conventions}. Collecting exponents gives the
lattice height $A$ and the number $V$ of unit-coset factors. The exponent is a
function of the lattice geometry of the cell alone; see \cite{denef1991functional}
for the cell-decomposition formalism. The explicit lattice exponent at arbitrary
OM order is recorded in the certificate.
\end{proof}

\begin{lemma}[\textsc{BB3}: residual count over a finite field]\label{lem:bb3}
Let $\bbF$ be a finite field with $|\bbF|=Q$. The number of monic non-squarefree
polynomials of degree $d\ge2$ over $\bbF$ is $Q^{\,d-1}$, independently of the
characteristic. Equivalently, the locus of residual polynomials that fail to be
separable --- the locus that triggers OM descent by Hypothesis~\textup{\ref{hyp:gmn}}(GMN3)
--- is the discriminant locus, of universal point count $Q^{\,d-1}$, with no
inseparable-irreducible exceptions in any characteristic.
\end{lemma}

\begin{proof}
The number of monic squarefree polynomials of degree $d\ge2$ over $\bbF$ is
$Q^{d}-Q^{d-1}$; hence the non-squarefree count is $Q^{d-1}$. Over a finite field
every irreducible polynomial is separable, because $\bbF$ is perfect: every
element is a $p$-th power, so $y^{p}-c=(y-c^{1/p})^{p}$ is never irreducible. Thus
``non-separable'' and ``non-squarefree'' coincide, and the descent-triggering
locus is the discriminant locus with the stated universal count.
\end{proof}

Lemma~\ref{lem:bb3} is the precise meaning of ``Frobenius is bijective on
points''. The wild inseparability lives entirely in the $p$-adic tree structure
--- the slopes and the lifting --- and never in the residue-field counts, which
are always the universal discriminant-stratum polynomials. This is what makes the
recursion characteristic-uniform.

% ============================================================================
\section{The cluster recursion}\label{sec:recursion}

We now prove that each per-shape cluster volume is a uniform rational function of
$q$. The argument is an induction on Okutsu--Montes order, resting only on
Lemma~\ref{lem:bb1} and Lemma~\ref{lem:bb3}.

\subsection{Splitting and the Vandermonde pushforward}\label{sec:split}

If $\bar f=\bar g\,\bar h$ with $\bar g,\bar h$ coprime mod $\fkm$, then Hensel's
lemma identifies $f\mapsto(g,h)$ with a measure-preserving product bijection: the
resultant $|\Res(g,h)|=1$ is a unit, so the conditional law of a cluster's
translated coefficients is Haar on the corresponding congruence cell. This is the
exact product structure that lets the recursion factor through coprime clusters.
On a fixed cluster, changing variables from coefficients to roots in the splitting
field introduces the Vandermonde Jacobian; on a stratum of fixed tree shape $T$
the pairwise valuations $v(\alpha_{i}-\alpha_{j})$ are the tree distances and are
constant, so
\[
  \prod_{i<j}|\alpha_{i}-\alpha_{j}|=q^{-d(T)}
\]
is a single $q$-power determined by the shape. The pushforward of Haar measure
over the stratum is then the residue-configuration count times $q^{-(\text{cell
codimension})}$, both explicit by Section~\ref{sec:counts}.

\subsection{Per-shape volume by induction on OM order}\label{sec:bb3inf}

\begin{proposition}[Per-shape uniform rationality]\label{prop:per-shape}
For each OM tree shape $T$ on $n$ roots, the
cluster volume $C_{T}(q)$ is a uniform rational function of $q$, the same at every
prime including wild $p\mid n$.
\end{proposition}

\begin{proof}[Proof outline]
Induct on the OM order $r$ of the shape. At order $r$, the cluster is partitioned
by the order-$r$ Newton polygon into cells; by Hypothesis~\ref{hyp:gmn}(GMN2)
these are lattice cells, so each has box volume a char-independent $q$-power by
Lemma~\ref{lem:bb1}. Within a cell, the order-$r$ residual polynomial lives over
the finite field $\bbF_{r}$ by Hypothesis~\ref{hyp:gmn}(GMN1); a multiplicity-$1$
residual factor is a leaf with explicit $(e,f)$ by (GMN3), contributing a
$\bbF_{r}$-point count that is a universal polynomial in $q$ by Lemma~\ref{lem:bb3},
and a multiplicity-$\ge2$ factor descends to a strictly smaller child shape, to
which the inductive hypothesis applies. The passage from this finite-field
\emph{count} to a Haar \emph{volume} at order $r$ is the order-$r$ analogue of
Lemma~\ref{lem:bb1}: by Theorem~\ref{hyp:gmn}(GMN1)--(GMN2) the order-$r$ residual
coordinates are free $\bbF_{r}$-digits over a lattice cell, so the per-cell volume
is the residual point count times the char-independent box volume of
Lemma~\ref{lem:bb1}, exactly as at order $1$. (At order $1$ this is elementary; the
order-$r$ form is a direct consequence of the Gu\`ardia--Montes--Nart coordinate
facts and the Igusa--Denef cell decomposition, and is independently
oracle-corroborated through order $4$; see Section~\ref{sec:provenance}.) The number
of cells, the slopes, and the residue degrees are all $p$-independent, and the
descent terminates at finite $p$-independent order by (GMN4). Finally, a single shape contains infinitely many
descent depths --- the totally interior rescaling $\{v(c_{i})\ge i\}$ loops the
recursion back to the global problem at a rescaled argument --- and these resum as
a geometric series with ratio a $q$-power, contributing a denominator
$1-q^{-(e(e+1)/2-1)}$ for a cluster of size $e$ (for $e=3$ this is $q^{5}-1$ up to
a unit, the classical $\Phi$-type denominator). The result is a rational function
of $q$ with cyclotomic-type poles, assembled from $p$-independent data; it is
therefore the same rational function at every prime.
\end{proof}

\begin{remark}[Where the wildness goes]\label{rem:wild-absorbed}
The only place a wild prime could intervene is in the residual polynomials at a
descent node, where in characteristic $p$ a repeated factor is inseparable
(Artin--Schreier). By Lemma~\ref{lem:bb3} this is invisible to the count: the
non-separable locus is the discriminant locus with the universal count $Q^{d-1}$,
the same geometric data as the tame ``repeated root'' locus at an odd prime, even
though the two loci are different subvarieties with different tangent cones. This
is the volume-side image of the referee's tangent-space obstruction: the loci that
the resolution cannot smooth uniformly have nonetheless the same volume in every
characteristic.
\end{remark}

% ============================================================================
\section{Assembling the densities}\label{sec:assemble}

\subsection{From cluster volumes to the monic density}\label{sec:linear-system}

By the tree decomposition \eqref{eq:tree-decomp} and Proposition~\ref{prop:per-shape},
the monic density $\alpha(n,\sigma;q)$ is a finite sum of cluster volumes, each a
uniform rational function of $q$. Multi-cluster types are linearized: the
coefficients relating the type densities to the cluster volumes form a finite
linear system over $\bbQ(q)$ whose entries are residue-configuration counts and
$q$-powers, all $p$-independent. Solving the system expresses $\alpha(n,\sigma;q)$
as a uniform rational function of $q$.

\subsection{The monic-to-projective bridge}\label{sec:bridge}

Both densities are the same stratification of binary degree-$n$ forms by their
residue divisor on $\bbP^{1}$, driven by the same intrinsic cluster laws
$L_{e}$:
\[
  \rho(n,\sigma;q)=\frac{1}{|\bbP^{n}(\bbF_{q})|}\sum_{D\in\Sym^{n}\bbP^{1}(\bbF_{q})}W_{D}[\sigma],
  \qquad
  \alpha(n,\sigma;q)=\frac{1}{q^{n}}\sum_{D\in\Sym^{n}\mathbb{A}^{1}(\bbF_{q})}W_{D}[\sigma],
\]
with the same weights $W_{D}$ built from the cluster laws. The only difference is
whether the point $[1:0]\in\bbP^{1}$ is admissible, i.e. $q+1$ projective versus
$q$ affine base points, together with the normalization $|\bbP^{n}(\bbF_{q})|$
versus $q^{n}$; the crux identity is $\Sym^{n}\bbP^{1}=\bbP^{n}$. Consequently the
projective density is the monic recursion read at the \emph{reciprocal} argument.
This is why the monic density is not palindromic --- indeed the monic
$\alpha(3,\{(1,1)^{3}\};2)$ is $4/93$, not symmetric under $q\mapsto1/q$ --- while
the $\bbP^{1}$-symmetric projective $\rho$ is palindromic. The bridge is itself
$p$-independent: it is built from point counts and is the volume-side incarnation
of the symmetry-for-$\alpha$ identity $\beta(e;q^{\delta})=\alpha(e;q^{-\delta})$.

\begin{corollary}[Uniform rationality]\label{cor:rationality}
For each $n$ and $\sigma$, the projective
density $\rho(n,\sigma;q)$ is a uniform rational function $R_{\sigma}(q)\in\bbQ(t)$,
the same at every prime including wild $p\mid n$. This is the rationality half of
Theorem~\textup{\ref{thm:main}}.
\end{corollary}

% ============================================================================
\section{The functional equation at wild primes}\label{sec:wild-symmetry}

It remains to prove palindromy, $R_{\sigma}(1/t)=R_{\sigma}(t)$. The argument is
the clean payoff of the decoupling: we do not prove the functional equation at
wild primes directly; we transfer it.

\begin{proposition}[Palindromy]\label{prop:palindromy}
The projective density satisfies $R_{\sigma}(1/t)=R_{\sigma}(t)$.
\end{proposition}

\begin{proof}
By Corollary~\ref{cor:rationality} a single rational function $R_{\sigma}\in\bbQ(t)$
computes the projective density at \emph{every} prime: $R_{\sigma}(q)=\rho(n,\sigma;q)$
for each prime power $q$, tame or wild. At a tame prime $p\nmid n!$ the tame functional
equation (the Igusa--Denef--Meuser symmetry and the tame factorization-density
literature, input~(iii) of Section~\ref{sec:standing}) furnishes a \emph{palindromic}
rational function $P_{\sigma}\in\bbQ(t)$, $P_{\sigma}(1/t)=P_{\sigma}(t)$, that also
computes the density there: $P_{\sigma}(q)=\rho(n,\sigma;q)$ for all but finitely many
tame prime powers $q$. Thus $R_{\sigma}$ and $P_{\sigma}$ agree at infinitely many
values --- for any fixed $n$ there are infinitely many tame primes, with unboundedly
large residue cardinalities --- so the polynomial numerator of $R_{\sigma}-P_{\sigma}$
has infinitely many roots and vanishes: $R_{\sigma}=P_{\sigma}$ in $\bbQ(t)$. Therefore
$R_{\sigma}(1/t)=P_{\sigma}(1/t)=P_{\sigma}(t)=R_{\sigma}(t)$ identically, in particular
at the wild primes $p\mid n$, where the tame functional equation itself is unavailable.
\end{proof}

Together, Corollary~\ref{cor:rationality} and Proposition~\ref{prop:palindromy}
prove Theorem~\ref{thm:main}.

% ============================================================================
\section{Examples}\label{sec:examples}

We record the small cases in closed form, with $\Phi=\Phi(q)=q^{4}+q^{3}+q^{2}+q+1$.
The densities below are projective, weight-normalized, and exact rational functions
of $q$; palindromy and the sum-to-one normalization are visible by inspection.

\begin{example}[Quadratics, $n=2$]\label{ex:n2}
The three types are split $\{(1,1),(1,1)\}$, inert $\{(1,2)\}$, and ramified
$\{(2,1)\}$:
\[
  \rho\bigl(2,\{(1,1)^{2}\};q\bigr)=\tfrac12,\quad
  \rho\bigl(2,\{(1,2)\};q\bigr)=\frac{q^{2}-q+1}{2(q^{2}+q+1)},\quad
  \rho\bigl(2,\{(2,1)\};q\bigr)=\frac{q}{q^{2}+q+1}.
\]
Each is invariant under $q\mapsto 1/q$ and the three sum to $1$. The ramified type
includes the wild prime $p=2$ with no exception. As an independent check, a direct
$2$-adic computation at the wild prime gives the \emph{monic} densities
$\alpha(2,\{(1,1)^2\},\{(1,2)\},\{(2,1)\};2)=(\tfrac13,\tfrac13,\tfrac13)$, exactly the
values at $q=2$ of the uniform rational functions assembled by the recursion --- the
wild prime $p=2$ obeys the same formula as every tame prime, with no exception. This is
genuinely subtle at $p=2$: a monic quadratic whose reduction is an inseparable square,
$x^2+ax+b\equiv(x+c)^2\pmod 2$, may be split, inert, \emph{or} ramified, so the
reduction no longer determines the type (Hensel's lemma fails at the double root), yet
the three densities still land on the uniform values. The \emph{projective} densities
$\rho$ at $q=2$, read off the formulas above, are instead $(\tfrac12,\tfrac3{14},\tfrac27)$:
the monic densities are uniform but not palindromic, whereas the $\bbP^{1}$-symmetric
projective densities are both, the monic-to-projective bridge of
Section~\ref{sec:bridge} accounting for the difference.
\end{example}

\begin{example}[Cubics, $n=3$]\label{ex:n3}
The five types are $\{(1,1)^{3}\}$, $\{(1,1),(1,2)\}$, $\{(1,3)\}$, the partially
ramified $\{(1,1),(2,1)\}$, and the totally ramified $\{(3,1)\}$:
\[
\begin{aligned}
  \rho\bigl(3,\{(1,1)^{3}\};q\bigr)&=\frac{(q^{2}+1)^{2}}{6\,\Phi}, &
  \rho\bigl(3,\{(1,1),(1,2)\};q\bigr)&=\frac{q^{4}+1}{2\,\Phi},\\[0.3em]
  \rho\bigl(3,\{(1,3)\};q\bigr)&=\frac{q^{4}-q^{2}+1}{3\,\Phi}, &
  \rho\bigl(3,\{(1,1),(2,1)\};q\bigr)&=\frac{q^{3}+q}{\Phi},\\[0.3em]
  \rho\bigl(3,\{(3,1)\};q\bigr)&=\frac{q^{2}}{\Phi}. & &
\end{aligned}
\]
All five are palindromic under $q\mapsto 1/q$ and sum to $1$. The totally ramified
type $\{(3,1)\}$ includes the wild prime $p=3$ with no exception --- precisely the
case $(n,p)=(3,3)$ at which the $S_{3}$-equivariant resolution fails, yet the
density is the same uniform rational function $q^{2}/\Phi$.
\end{example}

\begin{example}[The first wild atom, $n=4$, $p=2$]\label{ex:wild-atom}
The first genuinely wild phenomenon appears in a size-$4$ cluster with a single
Newton slope $\tfrac12$. Conditioning on a clean slope-$\tfrac12$ cluster, the
residual polynomial is a quadratic $R(y)=y^{2}+\bar b\,y+\bar a$ with $\bar a$
uniform in $\bbF_{q}^{\times}$ and $\bar b$ uniform in $\bbF_{q}$. The
descent-triggering locus --- $R$ has a repeated root --- is
\[
  \{\bar b^{2}=4\bar a\}\ \text{at odd }p
  \qquad\text{versus}\qquad
  \{\bar b=0\},\ R=(y+\sqrt{\bar a})^{2}\ \text{at }p=2,
\]
which is \emph{inseparable} at $p=2$ (Artin--Schreier). These are geometrically
different loci --- a punctured parabola versus a punctured line, with different
tangent cones --- but both are the image of the merge map $\rho\mapsto(y-\rho)^{2}$
and so have the \emph{same} measure $1/q$ in every characteristic, by
Lemma~\ref{lem:bb3}. Resolving the full slope-$\tfrac12$ cluster gives the
characteristic-uniform conditional cluster law
\[
  P\bigl((4,1)\mid\text{inseparable sub-stratum}\bigr)=\frac{q}{q+1},
\]
the wild $p=2$ degeneration yielding the \emph{same} $q/(q+1)$ as the tame double
point at an odd prime. This is the deepest instance of the mechanism: the
char-dependent geometry is invisible to the volume.
\end{example}

\begin{remark}[A deep-wild value, $n=6$]\label{rem:n6}
The recursion reaches the first deep-wild Okutsu--Montes order-$3$ branch at $n=6$:
a slope-$\tfrac12$ cluster whose order-$1$ residual is a cube $(y-\rho)^{3}$,
forcing a size-$3$ inseparable child at $p=2$. The conditional cluster law that
the recursion derives for the totally ramified outcome $(6,1)$ inside this branch
is
\[
  \beta_{(6,1)}(q)=\frac{q^{2}(q-1)(q^{2}+1)}{q^{5}-1},
  \qquad \beta_{(6,1)}(2)=\frac{20}{31},
\]
the same rational function at $p=2$ and at the odd primes. This value
$20/31\approx0.645$ is what a PARI \texttt{factorpadic} oracle confirms at every
prime tested, and it \emph{rejects} the naive guess $q/(q+1)=2/3$ (rejected at
roughly $11\sigma$ on a $5$-event $\chi^{2}$). We stress that $\beta_{(6,1)}$ is a
\emph{conditional cluster volume}, not the full projective type density
$\rho(6,\{(6,1)\};q)$; the latter is a more complicated rational function that is
itself palindromic. The cross-check is evidence for the recursion, not a proof;
see Section~\ref{sec:provenance}.
\end{remark}

% ============================================================================
\section{Provenance and verification}\label{sec:provenance}

\paragraph{Inputs and scope.}
Theorem~\ref{thm:main} is proved using two pieces of standard, published machinery,
cited throughout: the higher-order Okutsu--Montes algorithm of
Gu\`ardia--Montes--Nart \cite{gmn2012} and the Igusa--Denef $p$-adic cell
decomposition \cite{denef1991functional,igusa2000introduction}; the tame functional
equation \cite{del_corso_dvornicich_2000,john,caruso2022zeroes} enters once, to seed
the palindromy transfer of Section~\ref{sec:wild-symmetry}. For $n\le3$ the proof is
elementary and uses none of the higher-order Montes machinery: the wild Newton boxes
are empty (the Newton polygon alone decides the factorization, with no order-$\ge2$
residual invoked). The higher-order algorithm first engages at the order-$2$
Artin--Schreier atom of Example~\ref{ex:wild-atom} ($n=4,\ p=2$), where the argument
is additionally corroborated by an independent PARI oracle at the wild prime $p=2$.
We note for the record that the $r$-point pole structure and related finer questions
form a separate line of work, currently partial, and are not treated in this paper.

\paragraph{Cross-checks (evidence, not proof).}
The densities produced by the recursion have been compared, as exact rational
functions of $q$, against the Bhargava--Cremona--Fisher--Gaji\'c root-count
moments. The full Okutsu--Montes density engine reproduces the BCFG root-count
distribution $P(n,r;q)$ and all of its raw and factorial moments \emph{exactly}
(symbolic identity, $\mathtt{simplify}=0$) for every degree it reaches, $n\le 6$,
including the deep-wild order-$3$ type of Remark~\ref{rem:n6}, with the
exhaustive-enumeration gate $\sum_{\sigma}\rho(n,\sigma;q)=1$ holding symbolically.
A PARI \texttt{factorpadic} oracle confirms the individual type densities, with no
unpredicted types, at primes $2,3,5,7$ within Monte-Carlo noise. These corroborate
the result but are not part of the proof. The one step worth singling out is the
order-$r$ count-to-volume passage of Proposition~\ref{prop:per-shape} (the
residual-equidistribution step, ``M6'' in the working notes): it is elementary at
order~$1$ and follows at every order from the Gu\`ardia--Montes--Nart coordinate
facts together with the Igusa--Denef cell decomposition. We have additionally
checked it directly against the oracle through OM order~$4$.

\paragraph{The companion Lean formalization.}
A companion certificate accompanies this paper, documenting a Lean~4 /
\textsf{mathlib} formalization of the count-native path; we record here precisely
what it establishes. The algebraic skeleton of the proof is fully machine-checked:
the decomposition $\rho=\sum_{T}C_{T}$, the rationality of each cluster volume
$C_{T}$, the geometric self-loop collapse, the finite-field residual count, the
box volume, and the palindromy transfer are all genuine Lean theorems whose
\texttt{\#print axioms} footprints are the three Lean-core axioms
($\texttt{propext}$, $\texttt{Classical.choice}$, $\texttt{Quot.sound}$) only. The
capstone \texttt{goal\_theorem\_montes} assembles these into the statement that the
genuine counting density is a uniform rational and palindromic function of $q$,
with footprint $\{\texttt{propext},\texttt{Classical.choice},\texttt{Quot.sound},
\texttt{tame\_functionalEquation}\}$, conditional on a minimal bundle of hypotheses
carried as structure fields (not as additional axioms): the
Gu\`ardia--Montes--Nart count-shadow facts and a single box-Haar identification.

The one genuine limitation is specific to the \emph{formalization}, not to the
mathematics: \textsf{mathlib} has no $p$-adic Haar measure, so the identification of
the per-box $p$-adic volume with the separately proved lattice/finite-field value
cannot yet be constructed inside Lean and is carried as a hypothesis (the ``measure
wall''). For this reason the formalization should be described as \emph{verified
modulo the box-Haar identification}, rather than ``formally verified'' without
qualification. This is a statement about what mathlib can currently express; the
mathematical theorem itself is proved in the body of this paper.

\paragraph{Byline and collaboration.}
The byline of this manuscript is an editorial decision of the project lead,
A.~G. It credits Claude (Anthropic) as primary author of the present draft because
the proof architecture, the structural mechanism of Section~\ref{sec:significance}, the
worked examples, and the prose were produced in an AI-assisted development loop.
A.~G.\ contributed in a human advisory and verification-discipline role: setting
the research program, supplying the framing of the inputs used
(Section~\ref{sec:standing}), running and adjudicating the independent symbolic
and PARI-oracle cross-checks above, and enforcing the honesty discipline of this
section --- in particular that the cited machinery be named, that cross-checks be
labelled evidence rather than proof, and that the companion formalization be
described modulo the box-Haar identification it carries. We record this
division of labour transparently so that the provenance of every claim is clear.

% ---- BIBLIOGRAPHY ----
\bibliography{reference}

\end{document}
